\section{原子核B\:\color{red}\tiny 10.19-10.25}
\subsection{物理史}
\lsn{1}填写下面物理史内容：

\ding{172}\lunl{4em}提出了原子论的基本假设，认为物质是由微小的、不可分割的粒子组成的，这些粒子称为原子。

\ding{173}提出枣糕模型的物理学家是\underline{\hbox to 5em{}}，他发现了\lunl{5em}

\ding{174}\underline{\hbox to 4em{}}通过\(\alpha\)粒子散射实验提出了原子的\underline{\hbox to 5em{}}模型

\ding{175}\underline{\hbox to 5em{}}最先发现了天然放射现象

\ding{176}最先发现质子的物理学家是\underline{\hbox to 4em{}}，发现中子的物理学家是\underline{\hbox to 5em{}}

\ding{177}发现\(X\)射线的是\underline{\hbox to 3em{}}

\ding{178}第一次在实验室实现人工转变的物理学家是\underline{\hbox to 4em{}}

\lsn{2}填写关于三种放射线的内容：

\begin{table}[htbp]
	\begin{tblr}{c|p{5em}|p{5em}|p{5em}}
	\hline
	&\(\alpha\) 射线&\(\beta\)射线&\(\gamma\)射线\\
	\hline
	组成&&&\\
	\hline
    电荷量&&&\\
    \hline
	速度&&&\\
	\hline
	贯穿本领&&&\\
	\hline
	电离能力&&&\\
	\hline
	\end{tblr}
\end{table}

\subsection{四种原子核反应}
\lsn{3}式\ce{^{A}_{Z}X}中，A、Z、X的含义分别是什么
\vspace*{2cm}

\lsn{4}写出常见的粒子表示符号：

质子：\lunl{3em}\quad 中子:\lunl{3em}\quad 氢的三种同位素：\lunl{3em}、\lunl{3em}、\lunl{3em}

氦的两种同位素：\lunl{3em}、\lunl{3em}\quad 碳的两种同位素：\lunl{3em}、\lunl{3em}

正负电子\lunl{3em}、\lunl{3em}\quad 中微子\lunl{3em}

\lsn{5}所有的核反应方程都必须满足哪些守恒定律？
\vspace*{1.5cm}

\lsn{6}衰变反应：写出\(\alpha\)衰变一般性方程与本质方程：
\vspace*{1.5cm}

\newpage
\lsn{7}衰变反应：写出\(\beta\)衰变一般性方程与本质方程：
\vspace*{1.5cm}

\lsn{8}裂变反应：\ce{^{235}_{92}U}在慢中子轰击分裂为\ce{^{141}_{56}Ba}和\ce{^{92}_{36}Kr}，写出它的反应方程
\vspace*{1.5cm}

\lsn{9}聚变反应：下列方程均为聚变反应，则未知粒子\ce{X}是：
\begin{itemize}
    \item \ce{^2_1H + ^3_2H \to ^4_2He + X}\qquad \lunl{4em}
    \item \ce{^2_1H + ^2_1H \to X + ^1_1p}\qquad \lunl{4em}
    \item \ce{^2_1H + ^2_1H \to X + ^1_0n}\qquad \lunl{4em}
    \item \ce{^3_2He + ^3_2He \to ^4_2He + 2 X}\qquad \lunl{4em}
\end{itemize}

\lsn{11}人工核反应：写出下面两个典型的人工核反应方程

\(\alpha\)粒子轰击\ce{^{14}N}发现质子：\lunl{20em}

\(\alpha\)粒子轰击\ce{^{9}Be}发现中子：\lunl{20em}


\lsn{12}写出半衰期定义、公式
\vspace*{2cm}
\subsection{质能方程与结合能}
\lsn{13}什么是质量亏损？
\vspace*{2cm}

\lsn{14}怎样理解结合能与比结合能这两个概念？
\vspace*{4cm}

\lsn{15}阐述比结合能与原子核稳定性的关系
\newpage
\subsection{真题检测}
\lsn{16}\icode{N17-A-SI-L2-260930-009}
\nceeID[6]{2017新课标\Rmnum{1}}{4}
大科学工程“人造太阳”主要是将氘核聚变反应释放的能量用来发电，氘核聚变反应方程是：\ce{^2_1H + ^2_1H \to ^3_2He + ^1_0n}，已知\ce{^2_1H}的质量为\qty{2.0136}{u}，\ce{^3_2He}的质量为\qty{3.0150}{u}，\ce{^1_0n}的质量为\qty{1.0087}{u}，\qty{1}{u}=\qty{931}{MeV/c^2}。氘核聚变反应中释放的核能约为\paren[]
\begin{choices}
	\item \qty{3.7}{MeV}
	\item \qty{3.3}{MeV}
	\item \qty{2.7}{MeV}
	\item \qty{0.93}{MeV}
\end{choices}

\lsn{17}\icode{N23-A-SI-L2-260930-002}
\nceeID[6]{2023全国乙}{3}
2022年10月，全球众多天文设施观测到迄今最亮伽马射线暴，其中我国的“慧眼”卫星、“极目”空间望远镜等装置在该事件观测中作出重要贡献。由观测结果推断，该伽马射线暴在1分钟内释放的能量量级为\qty{e48}{J}。假设释放的能量来自于物质质量的减少，则每秒钟平均减少的质量量级为（光速为\qty{3e8}{m/s}）\paren[]
\begin{choices}
	\item \,\qty{e19}{kg}
	\item \,\qty{e24}{kg}
	\item \,\qty{e29}{kg}
	\item \,\qty{e34}{kg}
\end{choices}
\begin{answer}
	C\qquad 提示：\(E = mc^2\)
\end{answer}

\lsn{18}\icode{N22L-A-SI-L2-260930-001}
\nceeID[4]{2022辽宁}{2}
2022年1月，中国锦屏深地实验室发表了首个核天体物理研究实验成果。表明我国核天体物理研究已经跻身国际先进行列。实验中所用核反应方程为\ce{X + ^{25}_{12}Mg \to ^{26}_{13}Al},己知\(X\)、\ce{^{25}_{12}Mg}、\ce{^{26}_{13}Al}的质量分别为\(m_1\)、\(m_2\)、\(m_3\),真空中的光速为\(c\),该反应中释放的能量为\(E\)。下列说法正确的是\paren[]
\begin{choices}
	\item \(X\)为氘核\ce{^2_1H}
	\item \(E = (m_1 + m_2 + m_3)c^2\)
	\item \(X\)为氚核\ce{^3_1H}
	\item \(E = (m_1 + m_2 - m_3)c^2\)
\end{choices}
\begin{answer}
	D
\end{answer}

\lsn{19}\icode{N13H-A-SI-L2-260929-003}
\nceeID[2]{2013上海}{7}
在一个\ce{^{238}_{92}U}原子核衰变为一个\ce{^{206}_{82}Pb}原子核的过程中，发生\(\beta\)衰变的次数为\paren[]
\begin{choices}
	\item 6次
	\item 10次
	\item 22次
	\item 32次
\end{choices}

\lsn{20}\icode{N23Z-A-SI-L2-260627-001}
\nceeID[4]{2023浙江}{5}
“玉兔二号”装有核电池，不惧漫长寒冷的月夜，核电池将\ce{^{238}_{94}Pu}衰变释放的核能一部分转换成电能。\ce{^{238}_{94}Pu}的衰变方程为\ce{^{238}_{94}Pu \to ^4_2He + ^x_{92}U}，则\paren[]
\begin{choices}
	\item 衰变方程中的\(x\)等于233
	\item \ce{^4_2He}的穿透能力比\(\gamma\)射线强
	\item \ce{^{238}_{94}Pu}比\ce{^x_{92}U}的比结合能小
	\item 月夜的寒冷导致\ce{^{238}_{94}Pu}的半衰期变大
\end{choices}
\begin{answer}
	C\qquad 提示：比结合能越大越稳定
\end{answer}